\documentclass[12pt,a4paper]{article}
\usepackage[utf8]{inputenc}
\usepackage{amsmath,amssymb}
\usepackage{graphicx}

\title{Fiche de révision --- Démonstration du Drive}
\author{Publié depuis l'application}
\date{\today}

\newcommand{\R}{\mathbb{R}}
\newcommand{\dd}{\,\mathrm{d}}

\begin{document}
\maketitle
\tableofcontents

\section{Ce fichier est compilé dans le navigateur}
Ce document \textbf{LaTeX} est rendu directement dans l'application : sections
numérotées, listes, tableaux, équations, références croisées et macros
personnalisées comme $\dd$ ou $\R$.

\subsection{Une liste}
\begin{itemize}
  \item texte en \emph{italique} et en \texttt{chasse fixe} ;
  \item une formule inline $e^{i\pi} + 1 = 0$ ;
  \item une référence à l'équation~\eqref{eq:travail}.
\end{itemize}

\subsection{Des équations numérotées}
\begin{equation}
  \dd U = \delta W + \delta Q \label{eq:premier}
\end{equation}
\begin{align}
  \delta W &= -P_{ext}\dd V \label{eq:travail} \\
  W &= -\int_{V_1}^{V_2} P_{ext}\dd V \notag
\end{align}

\section{Un tableau}
\begin{table}[h]
\centering
\caption{Quelques grandeurs}
\begin{tabular}{|l|c|r|}
\hline
Grandeur & Symbole & Valeur \\
\hline
Vitesse de la lumière & $c$ & 299\,792\,458 m/s \\
Constante de Planck & $h$ & 6{,}626e-34 J.s \\
\hline
\end{tabular}
\end{table}

\begin{theorem}[Carnot]
Le rendement maximal d'un moteur est $\eta = 1 - \dfrac{T_0}{T}$.
\end{theorem}

\begin{proof}
Conséquence directe du second principe.
\end{proof}

Voir la figure~\ref{fig:demo} ci-dessous.
\begin{figure}[h]
\centering
\includegraphics[width=0.5\linewidth]{demo.png}
\caption{Image importée dans le Drive}
\label{fig:demo}
\end{figure}

\end{document}
